\makeatletter
\def\input@path{{../styles/}{../../styles/}{../../../styles/}{../}{../../}{../../../}}
\makeatother
\documentclass{ee122_notes}
\input{macros.tex}
\input{preamble.tex}
\renewcommand{\releasedate}{April 1, 2026}


\begin{document}
\section*{EE 122/ME 141 Week 10, Lecture 2: Examples of state-feedback control design}
\subsection*{Instructor: \instructor}
\subsection*{Date: \releasedate}


\textbf{Example 1.} Consider the first-order system
\[
\dot{x} = -0.5x + u, \qquad y = 2x.
\]

\begin{enumerate}[label=\textbf{Task \arabic*.}, leftmargin=*]

\item Write \(A\), \(B\), \(C\), and \(D\) matrices for the system.

\vspace{1.5cm}

\item Note the open-loop pole(s) of the system.

\vspace{0.5cm}

\item Physically, we want the closed-loop output \(y(t)\) to reach \(63\%\) of its final value in \(1\) second when a step input is applied. What is your desired pole?



\vspace{2.3cm}

\item Find the state-feedback controller $K$ and the feedforward gain \(k_f\)
\[
u=-Kx + k_f r
\]
that achieves your desired step response.

\vspace{1.8cm}

\end{enumerate}

\textbf{Example 2.} Consider the DC motor transfer function
\[
G(s)=\frac{2}{2s+1}.
\]

\begin{enumerate}[label=\textbf{Task \arabic*.}, leftmargin=*]

\item Write a state-space representation in reachable canonical form.

\vspace{2.4cm}

\item Note down \(A\), \(B\), \(C\), and \(D\).

\vspace{0.8cm}

\item Note the open-loop pole(s) of the system.

\vspace{0.8cm}

\item Physically, we want the closed-loop output \(y(t)\) to reach \(63\%\) of its final value in \(1\) second when a step input is applied. What should the desired pole be?

\textit{Hint:} For a first-order closed-loop system
\[
G_{cl}(s)=\frac{b}{s+a},
\]
a unit step input gives
\[
Y(s)=\frac{b}{s(s+a)},
\]
so
\[
y(t)=y(\infty)\left(1-e^{-at}\right).
\]

\vspace{2.3cm}

\item Find the state-feedback controller
\[
u=-Kx
\]
that achieves your desired pole.

\vspace{1.8cm}

\end{enumerate}

% \newpage

\textbf{Example 3.} Consider the second-order system in reachable canonical form
\[
\dot{x}=Ax+Bu,\qquad y=Cx+Du
\]
with
\[
A=\begin{bmatrix}
-2 & -13\\
1 & 0
\end{bmatrix},
\qquad
B=\begin{bmatrix}
1\\
0
\end{bmatrix},
\qquad
C=\begin{bmatrix}
0 & 1
\end{bmatrix},
\qquad
D=0.
\]

\begin{enumerate}[label=\textbf{Task \arabic*.}, leftmargin=*]

\item Run the open-loop step response in MATLAB:
\[
\texttt{step(ss(A,B,C,D))}
\]

What do you observe about the oscillation frequency of the system?

\vspace{2cm}

\item Note the open-loop poles of the system.

\textit{MATLAB hint:}
\[
\texttt{eig(A)}
\]

\vspace{1.5cm}

\item We want the closed-loop system to have a lower oscillation frequency than the open-loop system. Propose a pair of desired poles.

\vspace{2cm}

\item Use MATLAB to compute the state-feedback gain:
\[
\texttt{K = place(A,B,DP)}
\]
where \(\texttt{DP}\) is your vector of desired poles.

\vspace{1.5cm}

\item With state feedback \(u=-Kx\), write the closed-loop matrix
\[
A_{cl}=A-BK
\]
in terms of
\[
K=\begin{bmatrix}k_1 & k_2\end{bmatrix}.
\]

\vspace{2.2cm}

\item Compare the entries of \(A\) and \(A_{cl}\). Which coefficients change under state feedback? How does this connect to what we saw in controller canonical form in class?

\vspace{2.2cm}


\end{enumerate}

\end{document}